\chapter{Colour Blindness}
\index{Colour Blindness}

\section*{Definition and Overview}
Colour blindness, more accurately termed colour vision deficiency (CVD), is a reduced ability to perceive and distinguish certain colours. It is a sensory limitation of the visual system rather than a form of blindness: most affected people see the world normally in almost all respects, with visual acuity, depth perception, and low-light vision typically intact. The most common form makes it difficult to tell red from green; a much rarer form impairs the discrimination of blue from yellow; and a very small number of people have complete absence of colour perception (achromatopsia). Colour vision deficiency is usually inherited and present from birth, but it may also be acquired later in life as a result of disease, injury, medications, or ageing. Because the condition rarely causes discomfort or obvious visual impairment, many people live for years without realising they have it and may first learn of it during a routine eye test, school screening, or occupational health assessment.

\section*{Epidemiology}
Colour vision deficiency is the most common inherited disorder of vision in humans. Red--green deficiency affects approximately 1 in 12 men (about 8 per cent) and roughly 1 in 200 women (about 0.5 per cent), a striking sex difference that arises because the responsible genes are located on the X chromosome and men, who carry a single X chromosome, have no second copy to compensate for a faulty gene. Blue--yellow deficiency is far rarer, affects both sexes equally, and is usually due to an abnormality on an autosome. The condition occurs in all ethnic and racial groups, although the precise prevalence varies between populations; for example, red--green deficiency is somewhat less common in people of African ancestry and more common in some Northern European groups. Acquired colour vision problems are less common overall and tend to appear in middle and later life in association with diabetes, glaucoma, cataract, age-related macular degeneration, and various neurological and toxic exposures. Because most cases are mild, official statistics probably underestimate the true frequency of the disorder.

\section*{Aetiology and Pathophysiology}
Normal colour vision depends on the retina, which contains two principal classes of light-sensitive cells. Rods are highly sensitive photoreceptors that operate in dim light and convey no colour information. Cones are concentrated in the central retina (the macula and fovea) and mediate daylight, fine-detail, and colour vision. There are three populations of cones, each containing a photopigment tuned to respond preferentially to long (red, L-cones), middle (green, M-cones), or short (blue, S-cones) wavelengths, and the brain compares and combines their signals to reconstruct the full spectrum of perceived colour.

\begin{itemize}
    \item \textbf{Inherited (congenital) deficiency:} the most common cause, resulting from gene changes that alter or eliminate the red or green cone pigments. Because the L- and M-cone genes sit close together on the X chromosome, they are prone to recombination and mutation, which explains the high frequency of red--green deficiency. Blue--yellow deficiency follows an autosomal pattern.
    \item \textbf{Acquired deficiency:} damage to the cones, retina, optic nerve, or visual pathways from conditions such as diabetes, age-related macular degeneration, glaucoma, multiple sclerosis, Parkinson's disease, and chronic liver or kidney disease.
    \item \textbf{Toxic and drug effects:} certain medications (for example, some antimalarials, anti-arrhythmics, and erectile-dysfunction drugs) and industrial toxins such as organic solvents and heavy metals can impair colour discrimination.
    \item \textbf{Trauma:} direct injury to the eye or head injury affecting the occipital visual cortex.
    \item \textbf{Ageing:} progressive yellowing of the lens and gradual loss of cone function subtly alter colour perception in later life, particularly affecting the blue end of the spectrum.
\end{itemize}

In every form the fundamental problem is the same: the visual system no longer generates the normal pattern of cone signals that allows the brain to separate the affected colours, so hues that differ only in the affected dimension are perceived as identical.

\section*{Clinical Features}
The features of colour vision deficiency vary with its type and severity:

\begin{itemize}
    \item Difficulty distinguishing red from green, most often confusion between shades of red, green, brown, grey, and olive
    \item Trouble separating blue from yellow or violet in the rarer tritan forms
    \item Difficulty reading colour-coded information such as charts, maps, graphs, electrical wiring, and traffic signals when the lights are similar in brightness
    \item Problems identifying colours in dim light or when colours are muted or pastel
    \item Practical struggles such as matching clothes, selecting ripe fruit, reading meat doneness, or noticing sunburn, rashes, and bruising on the skin
    \item Characteristically, normal visual acuity, normal depth perception, and normal night vision in congenital forms, so the person does not complain of poor sight
    \item A history of difficulties at school or work with colour-based tasks that the person may have learned to compensate for silently
\end{itemize}

People with mild deficiency may be entirely unaware of their limitation and only become aware during formal testing. Complete achromatopsia, in contrast, is severe: affected individuals see the world in shades of grey and typically also have poor visual acuity, light sensitivity (photophobia), and nystagmus.

\section*{Differential Diagnosis}
Symptoms of colour confusion can be confused with, or coexist with, other eye problems:

\begin{itemize}
    \item \textbf{Cataract:} yellowing and darkening of the lens distorts and desaturates colour perception, particularly blues and purples
    \item \textbf{Glaucoma and optic nerve disease:} can impair colour vision early, before affecting other aspects of sight
    \item \textbf{Diabetic retinopathy and retinal disease:} damage to the macula or cones produces acquired colour deficits
    \item \textbf{Age-related macular degeneration:} destroys the central retina where cones are most dense
    \item \textbf{Optic neuritis:} an inflammatory demyelination of the optic nerve that classically reduces red--green discrimination
    \item \textbf{Medication and toxin effects:} reversible colour disturbance from certain drugs or chemical exposure
    \item \textbf{Normal age-related change:} a gradual, expected decline in colour discrimination in older people
\end{itemize}

The key clinical distinction is between congenital deficiency (present for life, stable, bilateral, and affecting both eyes equally) and acquired deficiency (new onset, potentially progressive, sometimes asymmetrical, and associated with other visual symptoms).

\section*{When to Seek Medical Care}
An optometrist or doctor should be consulted when:

\begin{itemize}
    \item Colour vision was previously normal and has become noticeably worse
    \item Colour changes appear suddenly or affect only one eye
    \item Colour problems are accompanied by blurred vision, eye pain, headaches, or floating spots
    \item An adult with diabetes, high blood pressure, or another chronic condition develops new colour difficulties
    \item A child seems unable to name or distinguish colours and you wish to arrange testing and support
    \item Colour deficiency appears alongside other neurological symptoms such as weakness or difficulty walking
\end{itemize}

\textbf{Contact your local emergency services for:}
\begin{itemize}
    \item Sudden loss of vision or a sudden change in vision in one or both eyes
    \item Sudden eye pain with redness and vision change together, which may indicate acute angle-closure glaucoma or other emergencies
    \item Sudden weakness, numbness, facial drooping, or difficulty speaking occurring with a change in vision, as these may indicate stroke
\end{itemize}

\section*{Diagnosis and Investigations}
Colour vision deficiency is diagnosed with simple, painless tests that can be performed during a routine eye examination:

\begin{itemize}
    \item \textbf{History:} asking about childhood colour difficulties, family history (a maternal uncle or grandfather with the condition), occupational needs, and use of medications
    \item \textbf{Ishihara plates:} the standard screening test, in which numbers or shapes are embedded in a field of coloured dots that people with red--green deficiency cannot read; most people with significant deficiency fail a set number of plates
    \item \textbf{Farnsworth--Munsell 100 hue and D-15 tests:} arranging coloured caps or chips in sequence to grade the type, axis, and severity of the deficiency
    \item \textbf{Anomaloscope:} an instrument that asks the person to match a mixture of red and green to a yellow field, providing a precise measure of red--green function and the type of abnormality
    \item \textbf{Full eye examination:} assessment of visual acuity, intraocular pressure, fundoscopy of the retina and optic disc, and a check of the pupils to identify acquired causes
    \item \textbf{Referral:} to an ophthalmologist, neurologist, or physician when an acquired, toxic, or systemic cause is suspected; additional imaging such as optical coherence tomography (OCT) may then be arranged
\end{itemize}

Screening of children can be undertaken informally at home and formally at school; children should be tested before starting school if there is any doubt, because colour-coded learning materials are widely used.

\section*{Management}
There is currently no cure for inherited colour vision deficiency, but much can be done to reduce its impact and to treat acquired forms:

\begin{itemize}
    \item \textbf{Adaptation and strategies:} learning to rely on position, brightness, pattern, and text rather than colour alone to identify and distinguish objects
    \item \textbf{Assistive technology:} smartphone applications and camera filters that identify and name colours, and that adjust the display for colour-blind users; screen overlays and operating-system colour filters on computers
    \item \textbf{Special lenses:} tinted glasses and contact lenses may enhance contrast for some wearers, though they do not restore normal colour vision and are not a cure; they can nevertheless be helpful in specific settings
    \item \textbf{Workplace and study accommodations:} requesting that colour-coded information be presented with patterns or labels in addition, and informing teachers, employers, and exam boards where relevant
    \item \textbf{Occupational guidance:} certain careers (aviation, rail signalling, electrical work, some military and maritime roles) have strict colour-vision requirements, and individuals should understand their options and any legal implications early
    \item \textbf{Treatment of the underlying cause:} in acquired cases, control of diabetes, treatment of glaucoma, cataract surgery, and removal of the offending medication may halt or reverse the colour deficit
    \item \textbf{Psychological support:} reassurance that the condition is not a sign of failing vision, and practical help for children to avoid the condition being mistaken for a learning difficulty
\end{itemize}

\section*{Complications}
Colour vision deficiency itself is benign and does not damage the eye, but it carries practical and occasionally safety-related consequences:

\begin{itemize}
    \item Safety hazards in some settings, including difficulty reading chemical hazard labels, electrical wiring, and signal lights
    \item Restricted access to certain occupations and roles with colour-vision standards
    \item Everyday difficulties with cooking, shopping, dressing, and interpreting charts and maps
    \item Frustration, embarrassment, and reduced confidence, particularly in childhood and adolescence
    \item Missed detection of medical signs, such as failing to notice jaundice, rashes, bruising, or a red, inflamed eye
    \item In acquired cases, progression of the underlying eye, nerve, or systemic disease if it remains untreated
\end{itemize}

\section*{Prognosis and Outcomes}
Inherited colour vision deficiency is a lifelong, stable condition that does not worsen over time and does not affect overall health or lifespan. Most people adapt fully and lead unrestricted lives, encountering only occasional practical difficulties. Children with the condition usually manage well at school once parents and teachers understand it, and most careers remain open to them. Acquired colour vision problems have a more variable outlook: they may improve if the underlying cause is treatable (for example, cataract extraction or drug withdrawal) but may progress if the underlying disease advances. Complete achromatopsia is a more serious condition with permanent visual impairment, but it remains non-progressive. The single most important factor in outcome for acquired cases is timely investigation and treatment of the cause.

\section*{Prevention and Self-Care}
\begin{itemize}
    \item Arrange vision testing for children, including school screening, so that colour vision deficiency is recognised early and not mistaken for a learning problem
    \item Keep up with eye reviews and treatment if you have diabetes, glaucoma, or another condition that can affect the retina or optic nerve
    \item Learn to use position, brightness, pattern, and labels rather than colour alone when identifying objects
    \item Add text labels and patterns to charts, maps, and diagrams rather than relying on colour coding
    \item Tell employers, teachers, and colleagues about your colour limitation wherever it is relevant to safety or performance
    \item Use a colour-identifying smartphone app for shopping, cooking, sorting laundry, and checking medication labels
    \item Wear eye protection during work and sport to reduce the risk of acquired colour loss from injury
    \item See an eye professional promptly if colour vision changes after previously being normal
\end{itemize}

\section*{Sources and Further Reading}
The following organisations provide reliable, accessible information on colour vision deficiency:

\begin{itemize}
    \item NHS. \textit{Colour vision deficiency (colour blindness) information.} https://www.nhs.uk/conditions/colour-vision-deficiency/
    \item MedlinePlus. \textit{Colour vision deficiency overview.} https://medlineplus.gov/colorvisiondeficiency.html
    \item National Eye Institute. \textit{Facts about colour blindness.} https://www.nei.nih.gov/
    \item MSD Manual. \textit{Colour blindness (colour vision deficiency).} https://www.msdmanuals.com/
    \item Mayo Clinic. \textit{Poor colour vision.} https://www.mayoclinic.org/
    \item World Health Organization. \textit{Blindness and vision impairment fact sheet.} https://www.who.int/
\end{itemize}
